\documentclass[12pt,a4paper]{article}
\usepackage[utf8]{inputenc}
\usepackage[portuguese]{babel}
\usepackage[T1]{fontenc}
\author{Ítalo César Alves de Oliveira}

\begin{document}
\begin{flushleft}
\textbf{\textsc{Introdução}}
\end{flushleft}

	Jogos digitais são atrativos por oferecer uma fuga momentânea da realidade. Com o avanço das tecnologias móveis cada vez mais dispositivos podem ser utilizados para executar jogos digitais, assim como os JBL (jogos baseados em localização) que necessitam de sensores como o GPS(Global Positioning System) e conexão com a internet.
	
	Os jogos baseados em localização são aqueles em que a posição geográfica de objetos (aparelho e outros elementos do jogo) é uma característica importante (ou faz diferença) para se jogar\cite{Luis}. Os Jogos Baseados em Localização (JBL) são jogos pervasivos que usam tecnologias  de localização para incorporar a posição física do jogador às suas regras e dinâmicas\cite{Nayana}. Em geral, um LBG pode ser jogado virtualmente em qualquer lugar, desde que seus elementos de jogo sejam mapeados corretamente para os locais apropriados no mundo real\cite{Maia}.
	
	Desenvolver LBGs em todo o mundo é uma tarefa desafiadora devido à necessidade de implantar instâncias de jogos em vários locais\cite{Maia}. Com estre trabalho se espera conseguir desenvolver um jogo JBL utilizando métodos, ferramentas e formulas para facilitar o desenvolvimento do mesmo, procurando qual ou quais serão mais uteis para o desenvolvimento do mesmo.
	
\begin{flushleft}
\textbf{\textsc{Explicação dos artigos:}}
\end{flushleft}

	Este artigo consiste em uma pesquisa para equilibrar os mapas de jogos baseados em localização por meio da transposição de mapas através de um algoritmo de transposição para combinar gráficos semelhantes entre outros para resolver problemas adicionais para atingir o equilíbrio\cite{Maia}.
	
	Este artigo produz uma pesquisa para averiguar se o conjunto de técnicas e ferramentas fornecidas pelo GUR (Games User Research), que serve para avaliar a experiencia proporcionada por jogos digitais e obter \textit{insights}, também se aplica aos JBL (Jogos baseados em localização)\cite{Nayana}.
	
	Este artigo efetuou uma analise com o intuito de determinar métodos para minimizar ou eliminar limitações tecnológicas que interferem no desenvolvimento de Jogos baseados em localização por meio de experimentos com esse métodos\cite{Luis}.
	
\begin{flushleft}
\textbf{\textsc{Resumo:}}
\end{flushleft}

	Todos os trabalhos citados possuem um bom desenvolvimento, o artigo \cite{Luis} possui potencial para ter impactos positivo para o desenvolvimento de JBL pois apesar de ser um artigo relativamente antigo ele aborda um ponto crucial para o desenvolvimento dos mesmos apesar das limitações, assim como o artigo citado anterior mente este artigo \cite{Nayana} analisa a possibilidade das técnicas e ferramentas fornecidas pelo GUR serem efetivas para os JBL oque auxiliaria na analise do impacto dos mesmos nos usuários, já o artigo \cite{Maia} possui mais foco em uma área que afeta mais o usuário que é o balanceamento do mapa sendo assim excelente para tornar os JBL atrativos para usuários independentemente de sua localidade
	
\bibliographystyle{plain}
\bibliography{ref}

\end{document}